\section{Introduction}
\label{sec:intro}
Iodine released from lead--bismuth eutectic (LBE) can be transported in
different chemical forms before depositing on cooler surfaces. The
chemical form matters because it determines both the gas-phase inventory
available for transport and the equilibrium pressure governing wall
exchange. Liu et al.\ \cite{Liu2025chemical} investigated this behaviour
by evaporating pure iodides and iodine-containing LBE into helium,
separating the products along a negative temperature gradient and
characterising the deposits. Their observations provide spatially
resolved evidence for PbI$_2$ and BiI$_3$ deposition, together with
potassium-containing deposits in some LBE experiments. Connecting such
measurements to a transport calculation requires more information than
a single deposition temperature: the temperature-coordinate mapping,
carrier flow, source history and gas composition all enter the predicted
distribution.

Lobresco et al.\ \cite{Lobresco2026assessment} compared temperature-based
and vapour-pressure-driven CFD deposition models. The former require
calibration of a deposition law; the latter connect deposition to
vapour-pressure information through interfacial kinetics. Their study
also motivates extending transport models to multiple chemical species
and thermodynamic equilibrium. A multispecies calculation introduces
an additional consistency requirement. Changes in chemical form must
preserve the transported elements, and different gas pathways leading
to the same condensed material must use one physical wall inventory.
An internally consistent implementation can nevertheless give uncertain
experimental predictions if the thermodynamic or source inputs are
poorly constrained.

The present work addresses these issues in the finite-volume framework
of OpenFOAM \cite{Weller1998tensorial}. Nine Pb--Bi--I gases are transported
in a prescribed helium carrier. Homogeneous gas equilibrium redistributes
the local elemental inventory, while separate interfacial channels
exchange material with immobile condensed wall reservoirs. The wall
flux includes both interfacial kinetic conductance and diffusion through
the distance from the neighbouring cell centre to the wall. This avoids
representing deposition by an arbitrarily chosen volumetric wall-layer
thickness. It does not remove spatial discretisation error.

Equilibrium-based reactive transport is also supported by general
thermodynamic packages such as GEM-Selektor and its GEMS3K kernel
\cite{Kulik2013gemselektor}. The software developed here contains a
GEMS3K interface, but the calculations reported in this paper use a
specialised ideal-gas element-potential solver with tabulated formation
constants. We therefore distinguish the thermodynamic formulation from
the particular numerical engine used in the reported simulations.
Condensed phases are handled through wall exchange rather than through
homogeneous cell equilibrium.

The study uses the published paper as its experimental input. No
author-supplied temperature files, release histories or mass-flow-controller
reference conditions are available. Instead of inferring a unique physical
source from deposition data, we report the consequences of two explicit
source constructions and selected variations in flow, diffusion,
interfacial accommodation and thermodynamic information. Separate
time-step and mesh studies identify which conclusions survive numerical
refinement. A temperature-coordinate discrepancy in the silica experiment
is examined as an input-consistency problem, with alternate mappings
defined from published quantities before comparing the resulting CFD
profiles.

The contribution is a documented conservative transport--equilibrium--wall
coupling and a completed numerical and input-sensitivity assessment.
The evidence comprises 15 baseline and 64 follow-up simulations of three
available experimental configurations. These are computational cases,
not independent experimental replicates. The scope is deliberately
restricted to the executed model: a fixed carrier and temperature field,
prescribed release, ideal-gas speciation, and wall deposits without aerosol
formation or surface-material-specific chemistry. Experimental comparisons
are conditional assessments of that model, rather than a claim of
independent predictive validation.
